\section{Supplementary Material}
\label{sec:appendix}

\subsection{Reproducibility Cookbook}
\label{sec:repro-cookbook}

The complete pipeline from a clean clone to a compiled PDF:

\begin{lstlisting}[language=bash, caption={Clean-clone-to-PDF transcript.}]
git clone https://github.com/auditidentity/blackrim-routing-caching-paper
cd blackrim-routing-caching-paper
pip install -r scripts/eval-routing/requirements.txt

python scripts/aggregate-cost-split.py \
  --telemetry /path/to/blackrim/.beads/telemetry/invocations.jsonl \
  --pricing /path/to/blackrim/internal/pricing/sheets/anthropic-public.toml \
  --session-evidence data/raw/session-evidence.json \
  > data/aggregated/cost-split.csv

python scripts/aggregate-cache-stats.py \
  --telemetry /path/to/blackrim/.beads/telemetry/invocations.jsonl \
  > data/aggregated/cache-stats.csv

python scripts/eval-routing/run.py \
  --labels scripts/eval-routing/labels.yml \
  > data/aggregated/routing-baseline-results.csv

python scripts/eval-plan-cache/index.py && \
  python scripts/eval-plan-cache/replay.py \
  > data/aggregated/plancache-summary.csv

make
\end{lstlisting}

The pre-aggregated CSVs shipped with this repository allow the LaTeX build
to proceed without access to the Blackrim production instance.
Running \texttt{make} from a clean clone with the pre-aggregated data
in \texttt{data/aggregated/} reproduces \texttt{main.pdf} exactly.

\subsection{Full Cost-Split Table}
\label{sec:full-cost-split}

\Cref{tab:cost-split-full} renders the complete cost decomposition
sourced from \texttt{data/aggregated/cost-split.csv}
(generated by \texttt{scripts/aggregate-cost-split.py}).
Dollar figures are computed against the
\texttt{internal/pricing/sheets/anthropic-public.toml} rate sheet
(claude-opus-4-7: \$5.00/MTok input, \$25.00/MTok output;
cache-read multiplier 0.10\(\times\); cache-creation 1.25\(\times\),
verified 2026-04-24).

\begin{longtable}{lrrrr}
\caption{Full \blackrim cost split (7-day window, 2026-05-09 through 2026-05-15).
  Source: \texttt{data/aggregated/cost-split.csv}.}
\label{tab:cost-split-full} \\
\toprule
Bucket & Cost (USD) & Share (\%) & Calls & Notes \\
\midrule
\endfirsthead
\multicolumn{5}{c}{\textit{(continued from previous page)}} \\
\toprule
Bucket & Cost (USD) & Share (\%) & Calls & Notes \\
\midrule
\endhead
\midrule
\multicolumn{5}{r}{\textit{(continued on next page)}} \\
\endfoot
\bottomrule
\endlastfoot
Session total     & 115.54 & 100.00 & 500  & All recorded calls \\
Main-thread (Gestalt) & 114.81 & 99.37 & 445 & Turns on main orchestrator \\
Subagent dispatch & 0.7338 & 0.64   & 55   & All spawned workers \\
Unclassified      & 0.0000 & 0.00   & 0    & No unattributed tokens \\
\end{longtable}

\paragraph{Pricing assumptions.}
The rate sheet in use at evaluation time:
\begin{itemize}
  \item claude-opus-4-7: \$5.00/MTok input, \$25.00/MTok output,
    \$6.25/MTok cache-creation (5-min), \$0.50/MTok cache-read.
  \item claude-sonnet-4-6: \$3.00/MTok input, \$15.00/MTok output,
    \$3.75/MTok cache-creation (5-min), \$0.30/MTok cache-read.
  \item claude-haiku-4-5: \$1.00/MTok input, \$5.00/MTok output,
    \$1.25/MTok cache-creation (5-min), \$0.10/MTok cache-read.
\end{itemize}
Future evaluations can reprice without re-running telemetry scripts by
updating the rate sheet (\texttt{internal/pricing/sheets/\allowbreak{}anthropic-public.toml})
and re-running \texttt{scripts/aggregate-cost-split.py}.

\subsection{Routing Test-Set Labels}
\label{sec:routing-test-labels}

\Cref{tab:routing-labels} lists the first 10 of the 50 hand-labelled
turns (49 non-ambiguous; turn-33 is labelled \emph{ambiguous} and
excluded from F1 computation).
The full label file is at
\texttt{scripts/eval-routing/labels.yml} in the repository.

\begin{longtable}{lllp{7cm}}
\caption{Routing test-set labels (first 10 of 50).
  Full set: \texttt{scripts/eval-routing/labels.yml}.
  Columns: turn ID, gold tier, semantic-similarity prediction,
  abbreviated rationale.}
\label{tab:routing-labels} \\
\toprule
Turn & Gold & Predicted & Rationale (abbreviated) \\
\midrule
\endfirsthead
\multicolumn{4}{c}{\textit{(continued)}} \\
\toprule
Turn & Gold & Predicted & Rationale \\
\midrule
\endhead
\bottomrule
\endlastfoot
turn-01 & opus   & sonnet & Architecture synthesis across design options (false-negative) \\
turn-02 & sonnet & sonnet & Implementation brief, well-specified \\
turn-03 & sonnet & sonnet & Docs/README update, no novel design \\
turn-04 & haiku  & sonnet & Deterministic git commit (false-positive: short text) \\
turn-05 & haiku  & haiku  & \texttt{/loop} slash-command invocation \\
turn-06 & opus   & sonnet & Multi-agent session-ID design (false-negative) \\
turn-07 & haiku  & haiku  & \texttt{/loop} slash-command invocation \\
turn-08 & opus   & sonnet & \texttt{/grill-me} architecture session (false-negative) \\
turn-09 & sonnet & sonnet & Orchestration wave status update \\
turn-10 & sonnet & sonnet & bd epic closure + follow-on dispatch \\
\end{longtable}

The dominant error pattern is opus-tier turns predicted as sonnet:
design-pivoting questions (\textit{``Do our agents need their own
session-ids?''}) are superficially similar to routine dispatch
turns in embedding space.
The 53-exemplar opus set is too small to form well-separated
clusters; this is the primary motivation for the v2 classifier
direction (\S\ref{sec:future}).

\subsection{Plan-Cache Signature Examples}
\label{sec:plancache-examples}

\Cref{tab:plancache-sigs} shows representative dispatch signatures
from the 40-dispatch synthetic fixture
(\texttt{tests/fixtures/\allowbreak{}sample-dispatches.jsonl}).
Each signature encodes the coarse tool-call sequence used to
key the plan cache (role, task shape, normalised file scope,
deliverable kind); file names are stripped to directory prefix.

\begin{longtable}{lp{10cm}}
\caption{Sample plan-cache signatures from the synthetic fixture.
  Source: \texttt{tests/fixtures/\allowbreak{}sample-dispatches.jsonl}.
  Notation: \texttt{Read:X/|Write:X/|Edit:X/|Bash} encodes the
  coarse ordered tool-call sequence.}
\label{tab:plancache-sigs} \\
\toprule
Task type & Signature \\
\midrule
\endfirsthead
\multicolumn{2}{c}{\textit{(continued)}} \\
\toprule
Task type & Signature \\
\midrule
\endhead
\bottomrule
\endlastfoot
build  & \footnotesize\texttt{Read:internal/|Bash:go|Edit:internal/} \\
build  & \footnotesize\texttt{Read:internal/|Write:internal/|Edit:internal/|Bash:go} \\
test   & \footnotesize\texttt{Read:internal/|Write:internal/|Bash:go} \\
build  & \footnotesize\texttt{Read:scripts/|Edit:scripts/|Bash:python} \\
prose  & \footnotesize\texttt{Read:mkdocs/|Write:mkdocs/} \\
read   & \footnotesize\texttt{Read:docs/|Read:internal/} \\
review & \footnotesize\texttt{Read:internal/|Read:internal/|Bash:grep} \\
operate & \footnotesize\texttt{Read:docs/|Read:infra/} \\
\end{longtable}

The six cache hits on the 40-dispatch fixture all occur on \texttt{build}
and \texttt{test} types with identical coarse tool-call sequences,
consistent with the expectation that implementation and test dispatches
are the most structurally repetitive.
The zero false-hit rate reflects the conservative exact-hash matching
on the coarse signature: two dispatches with the same sequence but
different deliverable kinds map to distinct keys.

\subsection{Literature Overview: Eight Surveyed Sub-Areas}
\label{sec:landscape}

The full Phase-A landscape research is at
\texttt{docs/research/tier-down-main-thread-landscape.md} in the \blackrim
repository. For reference, the eight sub-areas surveyed in that document
are listed below, ranked informally by relevance to the two primitives
presented in this paper.

\begin{enumerate}
  \item \textbf{Per-turn classification paradigms}
        (semantic similarity, DistilBERT-trained, LLM-as-router) ---
        direct precedent for §\ref{sec:routing}.
  \item \textbf{LLM routing}
        (RouteLLM~\cite{routellm}, Hybrid LLM Routing~\cite{hybridllm},
        AutoMix, Tryage) ---
        closest in problem framing; differ in targeting
        single-turn chat rather than long-horizon orchestration threads.
  \item \textbf{Model cascading}
        (FrugalGPT~\cite{frugalgpt} and successors) ---
        demonstrates that quality-preserving cost reduction via
        tier selection is achievable in production settings.
  \item \textbf{Confidence-based routing}
        (semantic entropy~\cite{semanticentropy}, hidden-state probes,
        verbalized confidence) ---
        natural successor to v1 router; requires live inference per turn.
  \item \textbf{Agentic plan caching}
        (\citeauthor{plancache}~\cite{plancache}) ---
        direct precedent for §\ref{sec:plancache}; measured on
        larger deployments with higher query repetition.
  \item \textbf{Cost-performance frontier benchmarks}
        (SWE-bench, BigCodeBench, Aider Polyglot) ---
        used indirectly to calibrate quality-preservation constraints.
  \item \textbf{Mixture-of-Experts at the orchestration layer} ---
        architectural alternative to per-turn routing;
        not applicable at the Claude Code API level.
  \item \textbf{Routing-eval methodology}
        (regret, expected utility, Pareto frontier) ---
        used to design the evaluation protocol in §\ref{sec:eval}.
\end{enumerate}

The cache-control deep-dive companion (commit 660054d on \blackrim main)
revised downward earlier estimates of CLAUDE.md reordering gains
(10--15\,\% incremental, not 43\,\%) and identified plan caching
as the higher-leverage opportunity at 3.2\,\% marginal savings on
the synthetic fixture (81.8\,\% composed vs.\ 78.6\,\% prefix-only).
